\documentclass[11pt,a4paper]{article}
\usepackage[utf8]{inputenc}
\usepackage[margin=1in]{geometry}
\usepackage{hyperref}
\usepackage{graphicx}
\usepackage{amsmath}
\usepackage{booktabs}

\title{RSCT Review: Judge Reliability Harness: Stress Testing the Reliability of LLM Judges}
\author{Swarm-It Research Discovery\\\small Automated RSCT-Based Analysis}
\date{March 06, 2026}

\begin{document}
\maketitle

\begin{abstract}
This document provides an automated review of the paper ``Judge Reliability Harness: Stress Testing the Reliability of LLM Judges'' from an RSCT (Representation-Solver Compatibility Testing) perspective. The paper achieved an RSCT relevance score of 6.2% and topic similarity of 44.2%.
\end{abstract}

\section{Paper Information}
\begin{itemize}
\item \textbf{Title:} Judge Reliability Harness: Stress Testing the Reliability of LLM Judges
\item \textbf{Authors:} Sunishchal Dev, Andrew Sloan, Joshua Kavner, Nicholas Kong, Morgan Sandler
\item \textbf{Source:} \url{https://arxiv.org/abs/2603.05399v1}
\item \textbf{RSCT Relevance:} 6.2%
\item \textbf{Topic Match:} 44.2%
\item \textbf{Key Concepts:} safety
\end{itemize}

\section{Original Abstract}
We present the Judge Reliability Harness, an open source library for constructing validation suites that test the reliability of LLM judges. As LLM based scoring is widely deployed in AI benchmarks, more tooling is needed to efficiently assess the reliability of these methods. Given a benchmark dataset and an LLM judge configuration, the harness generates reliability tests that evaluate both binary judgment accuracy and ordinal grading performance for free-response and agentic task formats. We evaluate four state-of-the-art judges across four benchmarks spanning safety, persuasion, misuse, and agentic behavior, and find meaningful variation in performance across models and perturbation types, highlighting opportunities to improve the robustness of LLM judges. No judge that we evaluated is uniformly reliable across benchmarks using our harness. For example, our preliminary experiments on judges revealed consistency issues as measured by accuracy in judging another LLM's ability to complete a task due to simple text formatting changes, paraphrasing, changes in verbosity, and flipping the ground truth label in LLM-produced responses. The code for this tool is available at: https://github.com/RANDCorporation/judge-reliability-harness

\section{RSCT Analysis}
This paper presents research with 6% relevance to RSCT theory.

\textbf{Abstract Summary:} We present the Judge Reliability Harness, an open source library for constructing validation suites that test the reliability of LLM judges. As LLM based scoring is widely deployed in AI benchmarks, more tooling is needed to efficiently assess the reliability of these methods. Given a benchmark dataset and an LLM judge configuration, the harness generates reliability tests that evaluate both binary judgment accuracy and ordinal grading performance for free-response and agentic task formats. We e...

\textbf{RSCT Relevance:} The work touches on concepts related to safety, which have connections to the Representation-Solver Compatibility Testing framework.

\textbf{Further Analysis:} A detailed manual review is recommended to fully assess the implications for RSCT-based certification approaches.


\subsection*{RSCT Certification Metrics}
\begin{tabular}{ll}
\textbf{Relevance (R):} & 0.375 \\
\textbf{Spurious (S):} & 0.318 \\
\textbf{Noise (N):} & 0.306 \\
\textbf{Kappa ($\kappa$):} & 0.550 \\
\end{tabular}

The kappa score of 0.550 indicates moderate representation-solver compatibility.


\section{Relevance to Swarm-It}
This paper was identified by the Swarm-It research discovery pipeline as potentially relevant to RSCT-based AI certification. The combined score of 21.4% places it in the exploratory category for further investigation.

\vspace{1em}
\hrule
\vspace{0.5em}
\small{Generated by Swarm-It Research Discovery | \url{https://swarms.network} | RSCT Certified}

\end{document}
